% !Mode:: "TeX:UTF-8"
\documentclass[12pt,a4paper]{article}
% \usepackage{ctex} % chinese support
\usepackage{amsmath,amscd,amsbsy,amssymb,latexsym,url,bm,amsthm}
\usepackage[vlined,boxed,ruled]{algorithm2e} %%不加end
\usepackage{epsfig,graphicx,subfigure}
\usepackage{enumitem,balance,mathtools}
\usepackage{wrapfig}
\usepackage{mathrsfs, euscript}
\usepackage[dvipsnames,table,xcdraw]{xcolor}
\usepackage{hyperref}
\usepackage{caption}

\newtheorem{theorem}{Theorem}
\newtheorem{lemma}{Lemma}
\newtheorem{proposition}{Proposition}
\newtheorem{corollary}{Corollary}
\newtheorem{exercise}{Exercise}
\newtheorem*{solution}{Solution}

\renewcommand{\thefootnote}{\fnsymbol{footnote}}

\newcommand{\postscript}[2]
{\setlength{\epsfxsize}{#2\hsize}
\centerline{\epsfbox{#1}}}

\renewcommand{\baselinestretch}{1.0}

\setlength{\oddsidemargin}{-0.365in}
\setlength{\evensidemargin}{-0.365in}
\setlength{\topmargin}{-0.3in}
\setlength{\headheight}{0in}
\setlength{\headsep}{0in}
\setlength{\textheight}{10.1in}
\setlength{\textwidth}{7in}
\makeatletter \renewenvironment{proof}[1][Proof]
{\par\pushQED{\qed}\normalfont\topsep6\p@\@plus6\p@\relax\trivlist
\item[\hskip\labelsep\bfseries#1\@addpunct{.}]\ignorespaces}{\popQED\endtrivlist\@endpefalse}
\makeatother
\makeatletter
\renewenvironment{solution}[1][Solution]
{\par\pushQED{\qed}\normalfont\topsep6\p@\@plus6\p@\relax\trivlist
\item[\hskip\labelsep\bfseries#1\@addpunct{.}]\ignorespaces}{\popQED\endtrivlist\@endpefalse}
\makeatother
\begin{document}
\noindent

%========================================================================
\noindent\framebox[\linewidth]{\shortstack[c]{
        \Large{\textbf{A Demo for Algorithm2e}}\vspace{1mm}\\
Help Document for CS2312-Algorithm and Complexity, Xiaofeng Gao@SJTU}}

\begin{enumerate}

    \item {\color{blue}\textbf{An example of If statements (with
        source code on the left and samples on the right):}}

        \begin{minipage}[t]{0.5\textwidth}

\begin{verbatim}
\begin{algorithm}[H]
    \LinesNumbered
    \KwIn{$x$, $y$}
    \KwOut{$sign$}
    \BlankLine
    \caption{$div(x,y)$} \label{Alg-div}
    \eIf{$rm(x,y)=0$}{
        $sign \leftarrow 1$\;
    }{
        $sign \leftarrow 0$\;
    }
    \Return{$sign$}\;
\end{algorithm}
\end{verbatim}
        \end{minipage}
        \hspace{2mm}
        \begin{minipage}[t]{0.4\textwidth}
            \begin{algorithm}[H]
                \LinesNumbered
                \KwIn{$x$, $y$}
                \KwOut{$sign$}
                \BlankLine
                \caption{$div(x,y)$} \label{Alg-div}
                \eIf{$rm(x,y)=0$}{
                    $sign \leftarrow 1$\;
                }{
                    $sign \leftarrow 0$\;
                }
                \Return{$sign$}\;
            \end{algorithm}
        \end{minipage}

    \item {\color{blue}\textbf{An example of If-ElseIf-Else statements:}}

        \begin{minipage}[t]{0.5\textwidth}
\begin{verbatim}
\begin{algorithm}[H]
    \LinesNumbered
    \KwIn{$score$}
    \KwOut{Letter Grade}
    \BlankLine
    \caption{LetterGrade($score$)}
    \label{Alg-Score}

    \uIf{$score \ge 90$}{
        \textbf{output} $A$\;
    }
    \uElseIf{$80 \le score < 90$}{
        \textbf{output} $B$\;
    }
    \Else{
        \textbf{output} $P$\;
    }
\end{algorithm}
\end{verbatim}
        \end{minipage}
        \hspace{2mm}
        \begin{minipage}[t]{0.4\textwidth}
            \begin{algorithm}[H]
                \LinesNumbered
                \KwIn{$score$}
                \KwOut{Letter Grade}
                \BlankLine
                \caption{LetterGrade($score$)} \label{Alg-Score}

                \uIf{$score \ge 90$}{
                    \textbf{output} $A$\;
                }
                \uElseIf{$80 \le score < 90$}{
                    \textbf{output} $B$\;
                }
                \Else{
                    \textbf{output} $P$\;
                }
            \end{algorithm}
        \end{minipage}

    \item {\color{blue}\textbf{An example of While statements:}}

        \begin{minipage}[t]{0.5\textwidth}
\begin{verbatim}
\begin{algorithm}[H]
    \LinesNumbered
    \KwIn{$x$, $y$}
    \KwOut{$x$}
    \BlankLine
    \caption{$rm(x,y)$} \label{Alg-rm}
    \While{$x \ge y$}{
        $x-=y$\;
    }
    \textbf{output} $x$\;
\end{algorithm}
\end{verbatim}
        \end{minipage}
        \hspace{2mm}
        \begin{minipage}[t]{0.4\textwidth}
            \begin{algorithm}[H]
                \LinesNumbered
                \KwIn{$x$, $y$}
                \KwOut{$x$}
                \BlankLine
                \caption{$rm(x,y)$} \label{Alg-rm}
                \While{$x \ge y$}{
                    $x-=y$\;
                }
                \textbf{output} $x$\;
            \end{algorithm}
        \end{minipage}

        \newpage
    \item {\color{blue}\textbf{An example of For statements:}}

        \begin{minipage}[t]{0.5\textwidth}
\begin{verbatim}
\begin{algorithm}[H]
    \LinesNumbered
    \KwIn{$n \in \mathbb{N}$}
    \KwOut{The sum from 1 to $n$}
    \BlankLine
    \caption{Sum($n$)} \label{Alg-Sum}
    $sum \leftarrow 0$\;
    \For{$temp=1$ to $n$}{
        $sum \leftarrow sum+temp$\;
    }
    \textbf{output} $sum$\;
\end{algorithm}
\end{verbatim}
        \end{minipage}
        \hspace{2mm}
        \begin{minipage}[t]{0.4\textwidth}
            \begin{algorithm}[H]
                \LinesNumbered
                \KwIn{$n \in \mathbb{N}$}
                \KwOut{The sum from 1 to $n$}
                \BlankLine
                \caption{Sum($n$)} \label{Alg-Sum}
                $sum \leftarrow 0$\;
                \For{$temp=1$ to $n$}{
                    $sum \leftarrow sum+temp$\;
                }
                \textbf{output} $sum$\;
            \end{algorithm}
        \end{minipage}

    \item {\color{blue}\textbf{An example of Repeat-Until statements:}}

        \begin{minipage}[t]{0.5\textwidth}
\begin{verbatim}
\begin{algorithm}[H]
    \LinesNumbered
    \KwIn{$a, b \in \mathbb{N}$}
    \KwOut{
        Greatest common divisor of $a$, $b$
    }
    \BlankLine
    \caption{GCD($a$, $b$)} \label{Alg-GCD}

    \Repeat{$gcd=0$}{
        $gcd \leftarrow a \mod b$\;
        $a \leftarrow b$\;
        $b \leftarrow gcd$\;
    }
    \textbf{output} $gcd$\;
\end{algorithm}
\end{verbatim}
        \end{minipage}
        \hspace{2mm}
        \begin{minipage}[t]{0.4\textwidth}
            \begin{algorithm}[H]
                \LinesNumbered
                \KwIn{$a, b \in \mathbb{N}$}
                \KwOut{Greatest common divisor of $a$, $b$}
                \BlankLine
                \caption{GCD($a$, $b$)} \label{Alg-GCD}

                \Repeat{$gcd=0$}{
                    $gcd \leftarrow a \mod b$\;
                    $a \leftarrow b$\;
                    $b \leftarrow gcd$\;
                }
                \textbf{output} $gcd$\;
            \end{algorithm}
        \end{minipage}

    \item {\color{blue}\textbf{An example of Case statements:}}

        \begin{minipage}[t]{0.5\textwidth}
\begin{verbatim}
\begin{algorithm}[H]
    \LinesNumbered
    \KwIn{$person$}
    \KwOut{$person$'s gender}
    \BlankLine
    \caption{Gender} \label{Alg-Gender}
    \Switch{$person$}{
        \uCase{$person.gender=male$}{
            \textbf{output} Male\;
        }
        \uCase{$person.gender=female$}{
            \textbf{output} Female\;
        }
        \Other{
            \textbf{output} Unknown\;
        }
    }
\end{algorithm}
\end{verbatim}
        \end{minipage}
        \hspace{2mm}
        \begin{minipage}[t]{0.4\textwidth}
            \begin{algorithm}[H]
                \LinesNumbered
                \KwIn{$person$}
                \KwOut{$person$'s gender}
                \BlankLine
                \caption{Gender} \label{Alg-Gender}
                \Switch{$person$}{
                    \uCase{$person.gender=male$}{
                        \textbf{output} Male\;
                    }
                    \uCase{$person.gender=female$}{
                        \textbf{output} Female\;
                    }
                    \Other{
                        \textbf{output} Unknown\;
                    }
                }
            \end{algorithm}
        \end{minipage}

    \item {\color{blue}\textbf{An example of complicated pseudocode:}}

        \begin{minipage}[t]{0.9\textwidth}
            \begin{algorithm}[H]
                \SetCommentSty{small}
                \LinesNumbered
                \KwIn{Alert set $\mathcal{A}$.}
                \KwOut{Flow rerouting set $\mathcal{M}_f^i$; VM migration
                set $\mathcal{M}_g^i$}
                \caption{Pre-Alert Management Procedure} \label{Alg_Check}
                $\mathcal{M}_f^i=\varnothing$,
                $\mathcal{M}_g^i=\varnothing$, $\textsc{Alert\_ToR}=\varnothing$\;
                \While{$\mathcal{A} \neq \varnothing$}{
                    Pick $\textsc{Alert}_x \in \mathcal{A}$\;
                    \Switch{$\textsc{Alert}_x$}{
                        \uCase{$\textsc{Alert}_x$ from $s_j$}{
                            {\color{ForestGreen}\tcp{ Compute congestion
                            alert from outer switch $s_j$}}
                            $\mathcal{F}=\{m_{ij}^k\mid \exists \mbox{
                                flows out from } m_{ij}^k \in h_{ij} \in v_i
                            \mbox{ passing through } s_j\}$\;
                            $\mathcal{M}_f^i = \mathcal{M}_f^i \cup
                            \textsc{Priority}(\mathcal{F}, \alpha)$\;
                        }
                        \uCase{$\textsc{Alert}_x$ from local ToR}{
                            $\textsc{Alert\_ToR}
                            =\textsc{Alert\_ToR}\bigcup \{\textsc{Alert}_x\}$\;
                        }
                        \Case{$\textsc{Alert}_x$ from $h_{ij}$}{
                            {\color{ForestGreen}\tcp{ Receive overload
                            alert from host $h_{ij} \in v_i$}}
                            $\mathcal{F}=\{m_{ij}^k\mid m_{ij}^k \in h_{ij}\}$\;
                            $\mathcal{M}_g^i = \mathcal{M}_g^i \cup
                            \textsc{Priority}(\mathcal{F},
                            1)$\; %\mbox{\textbf{Priority}}
                        }
                    }
                    $\mathcal{A}= \mathcal{A} \backslash \{\textsc{Alert}_x\}$\;
                }
                \If {$\textsc{Alert\_ToR}\neq \varnothing$} {
                    {\color{ForestGreen}\tcp{ Detect congestion alert
                    from $v_i$'s local ToR}}
                    $\mathcal{F}=\{m_{ij}^k\mid m_{ij}^k \in h_{ij},
                    \forall h_{ij} \in v_i \}$\
                    %Check whether flows from $m_{ij}^k \in h_{ij} \in v_i$;
                    $\mathcal{M}_g^i = \mathcal{M}_g^i \cup
                    \textsc{Priority}(\mathcal{F}, \beta)$\;
                }
                $\textsc{VmMigration}(\mathcal{M}_g^i)$
                {\color{ForestGreen} \tcp*[r]{Apply VM migration}} %
                % \mathcal{M}_g^i
            \end{algorithm}
        \end{minipage}

\end{enumerate}

{\color{blue} For more examples and explanations, please refer to the
AlgorithmPackage.pdf help document.}

%========================================================================
\end{document}

